\documentclass[11pt,a4paper]{article}

\usepackage[utf8]{inputenc}
\usepackage[T1]{fontenc}
\usepackage[french]{babel}
\usepackage{amsmath,amssymb,amsthm,mathtools}
\usepackage{geometry}
\usepackage{hyperref}
\geometry{margin=2.5cm}

\title{README --- Théorème A et vérification computationnelle}
\author{Mohamed Sofyen}
\date{\today}

\newtheorem{theorem}{Théorème}
\newtheorem{corollary}{Corollaire}
\newtheorem{proposition}{Proposition}

\begin{document}
\maketitle

\section{Objet}

Ce dépôt accompagne le travail sur les matrices de Cartier--Manin des
familles hyperelliptiques
\[
 y^2=x^d+\sigma x^r+t
\]
en caractéristique première impaire $p$.

Le résultat central présenté ici est une borne du rang obtenue à partir
du support arithmétique des coefficients et d'un couplage maximum dans
un graphe biparti.

\section{Théorème A}

On pose
\[
 n=\frac{p-1}{2},\qquad
 g=\left\lfloor\frac{d-1}{2}\right\rfloor
\]
et
\[
 f(x)=x^d+\sigma x^r+t.
\]
La matrice $M=M_{p,d,r}(\sigma,t)$ est la matrice $g\times g$ dont
l'entrée $(i,j)$ est le coefficient de $x^{ip-j}$ dans $f(x)^n$.

Pour $1\leq i,j\leq g$, on définit
\[
 T(i,j)=
 \left\{(\alpha,\beta,\gamma)\in\mathbb N^3:
 \alpha+\beta+\gamma=n,\quad
 d\alpha+r\beta=ip-j
 \right\}.
\]
Le graphe biparti $G_{p,d,r}$ possède les lignes et les colonnes comme
sommets et une arête $(i,j)$ lorsque $T(i,j)\neq\varnothing$.
On note $\nu(G)$ la taille d'un couplage maximum.

\begin{theorem}
Pour tous $\sigma,t$ dans un corps de caractéristique $p$,
\[
 \operatorname{rang}M_{p,d,r}(\sigma,t)
 \leq \nu(G_{p,d,r}).
\]
\end{theorem}

\subsection*{Idée de la preuve}

Le développement multinomial donne
\[
 M_{ij}
 =
 \sum_{(\alpha,\beta,\gamma)\in T(i,j)}
 {n\choose\alpha,\beta,\gamma}
 \sigma^\beta t^\gamma.
\]
Comme $n<p$, les coefficients multinomiaux sont non nuls modulo $p$.
Ainsi l'entrée est identiquement nulle exactement lorsque $T(i,j)$
est vide.

Tout mineur de taille strictement supérieure à $\nu(G)$ est nul :
un terme non nul de son développement de Leibniz fournirait un
couplage de cette taille, contradiction.

\section{Critère arithmétique des arêtes}

L'existence d'une arête $(i,j)$ équivaut à l'existence de
$\beta\in\{0,\ldots,n\}$ telle que
\[
 ip-j-r\beta\geq0,\qquad
 d\mid(ip-j-r\beta),
\]
et, en posant
\[
 \alpha=\frac{ip-j-r\beta}{d},
\]
on ait
\[
 \alpha+\beta\leq n.
\]
En particulier,
\[
 \gcd(d,r)\mid ip-j
\]
est une condition nécessaire.

Pour $d=p+1$, on retrouve la forme
\[
 ip-j=(i-1)d+(d-i-j),
\]
et donc la congruence
\[
 r\beta\equiv d-i-j\pmod d.
\]

\section{Conséquence de non-ordinarité}

Si
\[
 \nu(G_{p,d,r})<g,
\]
alors
\[
 \det M=0
\]
pour tous $\sigma,t$. Lorsque $f$ est sans facteur carré, cela entraîne
que la courbe hyperelliptique n'est pas ordinaire.

Attention : la condition $\nu(G)=g$ ne suffit pas à elle seule à prouver
l'ordinarité. De même, l'égalité
$\operatorname{rang}M=\nu(G)$ n'est pas vraie en général.

\section{Théorème B : une tranche arithmétique}

Pour
\[
 p\equiv5\pmod6,\qquad d=p+1,\qquad
 r=\frac{p+1}{3},
\]
le calcul combinatoire du graphe donne
\[
 \nu(G)=r=\frac{p+1}{3}
\]
et
\[
 g-\nu(G)=\frac{p-5}{6}.
\]
On obtient donc
\[
 \operatorname{rang}M\leq\frac{p+1}{3}
\]
et un défaut d'au moins
\[
 \frac{p-5}{6}
\]
par rapport au genre.

Le programme fourni vérifie cette structure pour des nombres premiers
choisis dans cette progression. Cette vérification numérique ne remplace
pas la démonstration mathématique du théorème.

\section{Ce que fait le programme}

Le fichier \texttt{theoreme\_couplage\_verif.py} :

\begin{enumerate}
\item construit les arêtes à partir de la définition exacte de $T(i,j)$ ;
\item calcule un couplage maximum par l'algorithme de Kuhn ;
\item construit la matrice de Cartier--Manin modulo $p$ ;
\item calcule son rang exact par élimination de Gauss sur $\mathbb F_p$ ;
\item vérifie numériquement
\[
\operatorname{rang}M\leq\nu(G).
\]
\end{enumerate}

Le code ne suppose jamais que le rang est égal au couplage maximum.

\section{Reproductibilité}

Avec Python 3 :
\begin{verbatim}
python3 theoreme_couplage_verif.py
\end{verbatim}

Aucune bibliothèque externe n'est nécessaire.

\section{Statut scientifique}

Le développement multinomial, la matrice de Cartier--Manin et le lien
entre p-rang et matrice de Cartier--Manin constituent le cadre classique.
Le travail présenté ici se concentre sur la construction arithmétique du
graphe $G_{p,d,r}$ et sur la borne de rang obtenue par son couplage
maximum.

Les résultats numériques doivent être distingués des résultats démontrés.
En particulier, les égalités de type
\[
\operatorname{rang}M=\nu(G)
\]
sont à traiter comme des observations ou des résultats à démontrer
séparément lorsqu'elles ne sont pas accompagnées d'une preuve.

\section{Références et positionnement}

Le document de référence mentionne notamment Miller, Lubin, Novoselov
et Zhu pour les cadres et familles proches. Ces travaux sont utilisés
pour le positionnement, non comme substitution à la preuve du présent
énoncé.

\end{document}
